\documentclass[a4paper,10pt]{article}
\newcommand\subdir{/home/koen/LaTeX-setup/}
\input{\subdir/LaTeX-files/imports.tex}
\input{\subdir/LaTeX-files/master-internship.tex}
\usepackage{\subdir/LaTeX-files/aastex}
\usepackage{natbib}
\bibliographystyle{\subdir/LaTeX-files/astroadstitle.bst}

%Set the title, Author and date
\title{Notes Week 23}
\author{Koen Essers}
\tutorial{Master Internship}
\date{\today}

%Set the header style
\header{\tutorialstring}

%Begin the document
\begin{document}
\setuplayout
\section{Grid}
Here I will note down my most important findings for the grid that is based on the accretion efficiency required to produce a \(1.29\;M_\odot \) barium star.
\begin{figure}[ht]
    \marginfignote{0}{This plot show the \(\varepsilon\) values for the different \(R_\text{Roche lobe}\) and \(q\) combinations. Note that the bottom right 3 squares are empty because they cannot produce a \(1.29\;M_\odot \) barium star, even with \(100\%\) accretion. The grid only goes up to \(q=0.8\) because \(q=0.9\) would produce barium stars heavier that \(1.29\;M_\odot \) even without any accretion.} 
    \centering
    \incplot{w22-eps-1000-small}
\end{figure}

Below I investigate the 3d\sidenote{\(5\times 5\times 5\)} grid that I simulated last week.

The first thing I want to check is \emph{how} close the final accretor masses are to the \(1.29\;M_\odot \) value.

If there would be no wind accretion after the initialization of the models, and they all interact at the predicted time, all models should have a mass of \(1.29\;M_\odot \).


\begin{figure}[ht]
    \marginfignote{0}{We produce many stars close to the \(1.29\;M_\odot \) final barium mass, but there is a spread with a maximum of about \(0.1\;M_\odot \). This looks quite bad here. Let's compare it with the actual barium star distribution.}
    \centering
    \incplot{w23-scatter}
\end{figure}

\begin{figure}[ht]
    \marginfignote{0}{Now we can see that the spread in the final masses of the \emph{MESA} models is quite a bit smaller than the spread of the Barium stars. }
    \centering
    \incplot{w23-scatter-compare}
\end{figure}



\begin{figure}[ht]
    \centering
    \incplot{w23-period}
\end{figure}

\begin{figure}[ht]
    \centering
    \incplot{w23-mass}
\end{figure}



\section{fix}

Some of the models in the grid interact directly after starting.
\begin{figure}[ht]
    \marginfignote{0}{This is not \emph{necessarily} wrong but it is \emph{better} to start the simulation a bit earlier, so we know for sure the stars would not have interacted at and earlier pulse.}
    \centering
    \incplot{w23-wrong-start}
\end{figure}


With the refactor of the old grid-generation code, I accidentally forgot to take the model BEFORE the first model that has a radius smaller than \(0.9\times R_\text{RL}\). Below, I show the \emph{new} and \emph{corrected} version that does have this feature again.
\begin{minted}[fontsize=\small,breaklines]{py}
def select_model(contact_age, models_dict):

    candidate = None
    models = list(models_dict.values())

    for i in range(len(models) - 1, -1, -1):

        if models[i]["age"] < contact_age:

            # Return the previous model if it exists
            if i > 0:
                return models[i - 1]

    print("no suitable stellar model before contact -> skipping")
    return candidate
\end{minted}

\newsection{Suitable MESA models}

The masses of the barium dwarfs range from \(0.73\) to \(1.5\;M_\odot \).
If we want to produce barium stars with these masses using MESA, and limit ourselves to a range of mass transfer efficiencies, we can exclude \emph{many} binary configurations. 

To see this in action, below I plot the mass of the final barium star after RLOF with an accretion efficiency of \(0.5\), which in this case, I assume is the \emph{upper} boundary of the efficiency\sidenote{It seems unlikely that for extreme RLOF like what is happening here, the efficiency will be higher than 0.5.}.
 

\begin{figure}[ht]
    \centering
    \incplot{w23-max-mass-region-all-masses}
\end{figure}

This means we can exclude all of the regions that are \emph{red} because, even with our highest amount of assumed accretion, they cannot produce stars with masses above \(0.73\;M_\odot \).

Similarly, the amount of transferred mass cannot be \emph{less than 0}. A natural lower bound for the accretion efficiency is then obviously \(\varepsilon =0\). Let's look at the masses of the barium stars that we produce without any RLOF accretion.

\begin{figure}[ht]
    \centering
    \incplot{w23-min-mass-region-all-masses}
\end{figure}

Now we can exclude stars that, even without any RLOF accretion, are heavier than our observed upper limit of \(1.5\;M_\odot \). These are the red regions in the figure above.

Together they form a limit on possible binary configurations. The plots show that this limit is mostly determined by the initial mass ratio, although the initial Roche lobe radius does alter the limits a little bit once winds start to be important.

The figures show another thing: namely the boundary between the systems that interact and undergo RLOF and systems that miss RLOF altogether. 

To illustrate this more clearly, I show the evolutionary paths for a constant \(q\) of \(0.8\).
\begin{figure}[ht]
    \centering
    \incplot{w23-wind}
\end{figure}
\nextpage
The green lines show the evolution of the Roche Lobe radius for systems that \emph{do not} undergo RLOF. The red lines show systems that \emph{do} undergo RLOF.

Since the most interesting parts seem to occur towards the end, I show the same figure again, but now zoomed in on last few pulses.
\begin{figure}[ht]
    \centering
    \incplot{w23-wind-zoom}
\end{figure}

For some masses, even a relatively equal mass ratio of \(0.8\) can cause RLOF interactions \emph{after} the last pulse. This depends on the mass of the convective envelope during the last pulse.

For the \(1.7\) and \(1.9\;M_\odot \) TPAGB star, we see a relatively large final bump, while the \(1.6\) and \(1.8\;M_\odot \) stars for example, have a much less pronounced final bump.

\begin{figure}[ht]
    \centering
    \incplot{w23-bump}
\end{figure}

The figure above shows exactly this. It is \emph{not} an effect of initial mass. It is really based on the mass left during the last pulse.

\nextpage

\begin{figure}[ht]
    \centering
    \incplot{w23-CO-region-all-masses}
\end{figure}

\begin{figure}[ht]
    \centering
    \incplot{w23-m_DUP-region-all-masses}
\end{figure}

\bibliography{master.bib}
\end{document}
